\scanpage{47}\mspage{$\overline{\underline{\mathrm{XVIII}}}$}
mais $\Lambda_{zy}\cap{]s,s'[}=\phi$ puisque $\Lambda_{zy}\subset\mathrm{int}\,\Lambda^1$, et par la même raison $\Lambda_{zy}\cap(\Lambda^2-\{s,s'\})=\phi$ car $\Lambda^2-\{s,s'\}\subset\mathrm{ext}\,\Lambda^1$ (cf. lemme $\overline{\underline{\mathrm{XVIII}}}$) cela prouve la première observation, on obtient la seconde en changeant $\Lambda^1$ par $\Lambda^2$ dans la preuve ci-dessus, c.q.f.d.

\uline{Lemme} $\overline{\underline{\mathrm{XX}}}$. Il y'a deux bons sommets de $\Gamma$, \inserted{non consécutifs si $\unsure{n|p}=$ nb de \unsure{cotés|côtés} de $\Gamma\geq4$,} \deleted{au moins} tels que les \inserted{interieurs des} triangles corespondants soient contenus dans $\mathrm{int}\,\Gamma$.

\uline{Démonstration} : L'hypothèse de récurrence appliquée à $\Lambda^1$ et $\Lambda^2$ montre qu'il'y a un \uline{bon sommet}\margin{tel que\ldots} % slanted
$t_1$ dans $\Lambda^1$ et un bon sommet $t_2$ dans $\Lambda^2$. Si $t_1\ne s,s'$ alors $t_1$ est un bon sommet de $\Gamma$ puisque $\mathrm{int}\,\Lambda^1\subset\mathrm{int}\,\Lambda$. (tout point de $\mathrm{int}\,\Lambda^1$ peut être relié à un point de $\mathrm{int}\,\Lambda$, par example en prenant un disque $D_x$, $x\in\Gamma^{1}_{ss'}-\{s,s'\}$ vérifiant b) de la def.1 par rapport à $\Gamma$ et à $\Lambda^1$ on voit que $\mathrm{int}\,\Lambda^1\cap\mathrm{int}\,\Lambda\ne\phi$....etc)

Dans le cas contraire sans perte de généralité on peut poser $t_1=s$. Deux cas \margin{?}peuvent se produire $ {]\text{\textcircled{\scriptsize \(s''\)}},s'[}\cap\Gamma=\phi$ ou bien $ {]s'',s'[}\cap\Gamma\ne\phi$ et dans ce cas $ {]s'',s'[}\cap S(\Gamma)\ne\phi$ puisque $ {]s'',s'[}\not\subset\Gamma$.

\margin{\deleted{\unsure{faux}}} % slanted
Si $\deleted{\unsure{[}}{]s'',s'[}\cap\Gamma=\phi$ il est clair que $s''$ est visible de $s'$ par l'interieur (car le triangle de sommets $s,s',s''\subset\mathrm{int}\,\Gamma$). Considerons alors $\unsure{\Lambda^{1'}_{\mkern5mu\cdot}|\Lambda^{1!}|\Lambda^{11}_{\mkern5mu\cdot}}$ le polygone \margin{!}obtenu en récollant selon \deleted{\illegible{1}} $\{s',s''\}$ l'arc de polygone $[s',s'']$ avec le seul des arcs de polygone de $\Lambda^1$ réliant $s$ à $s'$ qui ne contient pas $[s,s']$. Il est clair que $S(\Lambda^{1'})\subset S(\Gamma)$ et que l'ensemble de toutes les arêtes de $\Lambda^{1'}$ sauf une, savoir $[s',s'']$ est contenu dans $A(\Gamma)$, Mais nécessairement \margin{?}l'ensemble des arètes de $\Lambda^{1'}$ contenues dans $A(\Gamma)$ est de cardinal inferieur \margin{so what} % slanted
au cardinal de l'ensemble des arêtes de $\Lambda^1$ contenues dans $A(\Gamma)$.

Si $ {]s'',s'[}\cap S(\Gamma)\ne\phi$ on choisit $z\in S(\Gamma)\cap{]s'',s'[}$, il est clair que $z$ est visible \inserted{de $s=t$} par l'interieur (puisque $ {]s,z[}\subset\overset{\circ}{\text{triangle}}$ de sommets $s,s',s''$). à partir du couple $s,z$ on peut appliquer la même construction que ci-dessus pour aboutir aussi à un polygone $\Lambda^{1'}$ tel que : a) $S(\Lambda^{1'})\subset S(\Gamma)$; b)
% archive sheet number: 46
